% Computer-aided design (CAD) is a fundamental component of modern engineering, providing editable and parametric representations for downstream simulation, optimization, and manufacturing. Learning-based CAD generation has progressively moved from modeling CAD construction sequences and design histories~\citep{wu2021deepcad,willis2021fusion} toward language-driven generation with large language models (LLMs)~\citep{zhang2026large}. Recent methods can translate natural-language descriptions into parametric CAD sequences~\citep{khan2024text2cad}, generate CAD programs from textual or multimodal inputs~\citep{li2025llm4cad,sun2025large}, or directly synthesize executable modeling code through iterative or tool-augmented workflows~\citep{badagabettu2024query2cad,xie2025text,mallis2025cad}. In particular, CadQuery provides an executable, Python-based parametric representation that aligns naturally with the code-generation capabilities of pretrained LLMs~\citep{xie2025text,guan2026cad}.

% Despite this progress, reliably translating textual specifications into precise engineering structures remains challenging. Existing Text-to-CAD systems are commonly optimized or evaluated using geometric similarity, or visual consistency~\citep{khan2024text2cad,alrashedy2025generating,xie2025text,guan2026cad}. However, natural-language descriptions do not always establish an exact correspondence with the dimensions, topology, and structural details required by an engineering design, and small geometric deviations may substantially affect the behavior of the resulting structure. More importantly, engineering design is ultimately concerned with whether a component satisfies prescribed requirements under specific materials, loads, boundary conditions, and dimensional constraints, rather than merely whether it resembles a target shape. In conventional workflows, bringing an initially generated model into compliance often requires repeated manual measurement, geometric modification, simulation, and re-evaluation, resulting in a costly and labor-intensive refinement process. These limitations have motivated increasing attention to functional and physics-based criteria for evaluating and refining generated CAD models, allowing engineering feedback to guide geometry toward designs that satisfy downstream requirements~\citep{berger2026physics,son2026self}.

% Computer-aided engineering (CAE), particularly the finite element method (FEA), provides a natural source of deterministic feedback for evaluating generated designs. Recent studies have demonstrated the use of LLMs in automated FEA workflows and the acquisition of verifiable engineering signals~\citep{hou2025autofea,qi2025feagpt}. Meanwhile, reinforcement learning (RL) provides a natural mechanism for exploiting such automatically verifiable feedback~\citep{guan2026cad}, enabling simulation outcomes to be further incorporated into iterative CAD generation and optimization~\citep{berger2026physics,son2026self,deng2026cosmo}. Together, these approaches allow generated designs to be evaluated from both geometric and engineering perspectives. However, the resulting verification objectives are not independent, but exhibit a clear prerequisite structure: basic program validity and geometric compliance determine whether a generated design can be reliably carried forward to downstream engineering analysis, while physics-based evaluation becomes meaningful only after these lower-level requirements have been satisfied. Simply combining all objectives into a single flat reward may obscure these dependencies and make it difficult to effectively optimize objectives at different levels simultaneously. This motivates a hierarchical training strategy that progressively incorporates verification objectives according to their underlying dependency structure.

% In this work, we formulate physics-constrained Text-to-CAD generation, where a model receives textual geometric specifications together with materials, loading and boundary conditions, and physics-based constraints, and generates executable CadQuery programs. We construct a dataset for this setting and develop a complete CAD--CAE verification pipeline that automatically executes generated programs, verifies geometric and dimensional requirements, performs finite-element simulation, and converts deterministic verification outcomes into targeted reward signals. Building on supervised fine-tuning, we further introduce a hierarchical reinforcement-learning framework based on group-relative policy optimization~\citep{shao2024deepseekmath}, in which verification objectives are progressively incorporated according to their dependency structure. This enables physics-based verification to directly participate in model optimization, providing a controlled setting for studying how geometric and physical feedback can jointly guide Text-to-CAD generation.

CAD is a fundamental component of modern engineering, providing editable and parametric representations for downstream simulation, optimization, and manufacturing. Learning-based CAD generation has progressively moved from modeling CAD construction sequences and design histories~\citep{wu2021deepcad,willis2021fusion} toward language-driven generation with large language models (LLMs)~\citep{zhang2026large}. Recent methods can translate natural-language descriptions into parametric CAD sequences~\citep{khan2024text2cad}, generate CAD programs from textual or multimodal inputs~\citep{li2025llm4cad,sun2025large}, or synthesize executable modeling code through iterative or tool-augmented workflows~\citep{badagabettu2024query2cad,xie2025text,mallis2025cad}. In particular, CadQuery provides an executable, Python-based parametric representation that aligns naturally with code-generation capabilities of pretrained LLMs~\citep{xie2025text,guan2026cad}.

Despite this progress, geometric correctness alone does not ensure physical feasibility. Existing Text-to-CAD systems are commonly optimized or evaluated using geometric similarity or visual consistency~\citep{khan2024text2cad,alrashedy2025generating,xie2025text,guan2026cad}. 
However, natural-language descriptions do not always fully specify the dimensions, topology, and structural details required for engineering design, and even small geometric deviations may substantially affect structural behavior. More importantly, engineering design is ultimately concerned with whether a component meets prescribed dimensional and physical requirements under specified material properties, loads, and boundary conditions, rather than merely whether it resembles a target shape.
Bringing an initially generated model into compliance often requires repeated manual measurement, geometric modification, simulation, and re-evaluation, resulting in a costly refinement process. These limitations have motivated increasing attention to functional and physics-based criteria for evaluating and refining generated CAD models, allowing engineering feedback to guide geometry toward designs that satisfy downstream requirements~\citep{berger2026physics,son2026self}.

Computer-aided engineering (CAE), particularly simulation based on the finite element method (FEM), provides a natural source of deterministic feedback for evaluating generated designs. Recent studies have demonstrated the use of LLMs in automated FEM-based workflows and the acquisition of verifiable engineering signals~\citep{hou2025autofea,qi2025feagpt}. Simulation outcomes have also been used to guide iterative CAD generation and optimization~\citep{berger2026physics,son2026self,deng2026cosmo}. Reinforcement learning (RL) has further enabled CAD-generation models to learn from geometric feedback~\citep{guan2026cad}. Training with both geometric and physical feedback introduces a key challenge: prioritizing joint constraint satisfaction while retaining informative feedback on partial progress. A flat aggregation of heterogeneous rewards can obscure this priority, allowing improvements in individual criteria to compensate for unmet design requirements. This calls for a learning signal that distinguishes progress toward feasibility from improvements within the feasible design space.

As illustrated in Fig.~\ref{fig:teaser}, we present PhyCAD, a physics-aware Text-to-CAD framework that generates executable CadQuery programs from geometric specifications and physical constraints under prescribed material properties, loads, and boundary conditions. The task defines a constrained design space with admissible dimensional ranges rather than a unique target geometry. We construct a corresponding dataset and an automated CAD--CAE verifier that provides training feedback through program execution, geometric checks, and FEM-based simulation. Building on supervised fine-tuning, we introduce a Hierarchical Balanced Advantage for Group Relative Policy Optimization (GRPO)~\citep{shao2024deepseekmath}, grouping candidates by verification status and combining between-level progress with objective-wise peer ranking within each level. This formulation preserves cross-level priority in the training signal while providing fine-grained feedback for satisfying constraints and improving feasible designs. The trained model directly generates programs without an iterative simulation-guided refinement loop at inference time. Experiments across multiple language models show that PhyCAD consistently improves joint geometric and physical constraint satisfaction.


Our contributions are summarized as follows:
\begin{itemize}
[leftmargin=*, itemsep=0.15em, topsep=0.1em, parsep=0pt, partopsep=0pt]
\item We formulate physics-aware Text-to-CAD generation and construct a dataset that jointly captures geometric specifications, loading and boundary conditions, and physics-based constraints.
\item We develop an automated CAD--CAE verification pipeline that provides verifiable feedback on program executability, geometric and dimensional compliance, and physical constraint satisfaction to support model training.
\item We introduce a Hierarchical Balanced Advantage that combines between-level progress with objective-wise peer ranking within each level, preserving cross-level priority in the training signal while providing fine-grained feedback for satisfying constraints and improving feasible designs.
\end{itemize}



% Despite this progress, geometric correctness alone is insufficient for engineering design. Existing Text-to-CAD systems are commonly optimized or evaluated using program validity, geometric similarity, or visual consistency~\citep{khan2024text2cad,alrashedy2025generating,xie2025text,guan2026cad}. These criteria measure whether a generated program executes or reproduces the intended shape, but do not directly determine whether the resulting artifact satisfies engineering requirements under prescribed materials, loads, boundary conditions, and dimensional constraints. A geometrically plausible design may still exhibit excessive displacement, high stress, insufficient stiffness, or invalid load transfer. Recent studies therefore argue for evaluating generative CAD using functional and physics-based criteria rather than geometry alone~\citep{berger2026physics,son2026self}.

% This has motivated the integration of computer-aided engineering (CAE), particularly the finite element method (FEA), into LLM-based engineering workflows. AutoFEA and FeaGPT demonstrate that LLMs can assist or orchestrate FEA configuration, geometry processing, meshing, simulation, and result analysis~\citep{hou2025autofea,qi2025feagpt}. Physics-in-the-loop CAD agents further use deterministic simulation feedback to iteratively revise generated designs~\citep{berger2026physics,son2026self}, while COSMO-Agent formulates closed-loop CAD--CAE optimization as a tool-augmented reinforcement-learning problem with multi-constraint rewards~\citep{deng2026cosmo}. Meanwhile, CAD-Coder shows that reinforcement learning can improve Text-to-CAD generation using geometric and format rewards~\citep{guan2026cad}. These advances establish the value of verifiable feedback, but leave open how to organize heterogeneous CAD--CAE signals when they exhibit clear prerequisite dependencies: a candidate must first be executable, then geometrically compliant, before its physical performance can be meaningfully evaluated.

